\clearpage
\VerdierPhysicalPageAnchor{0147}%
\begingroup
\small
\setlength{\parskip}{0pt}%
\setlength{\emergencystretch}{1.5em}%
\setlength{\abovedisplayskip}{3pt plus 1pt minus 1pt}%
\setlength{\belowdisplayskip}{3pt plus 1pt minus 1pt}%

\VerdierSemanticAnchor{ch2-2-3-4-d-complete}
\noindent Ceci démontre d) ; démontrons maintenant f).

\VerdierSemanticAnchor{ch2-2-3-4-f}
Un objet $Z$ de $\mathcal D$ appartient à
$Q^{-1}(\operatorname{Lib}_R(Q))$ si et seulement s'il existe un objet
$Q$-libre à droite $Y$ et un isomorphisme dans $\mathcal D/\mathcal B$ :
$Q(Z)\xrightarrow[\sim]{m}Q(Y)$.

\VerdierSemanticAnchor{ch2-2-3-4-f-fraction}
Un tel isomorphisme est de la forme $m=Q(s)^{-1}Q(t)$, où
$s\in S_{\mathcal D}(\mathcal B)$ et
$t\in\widehat{S_{\mathcal D}(\mathcal B)}$ ((2.2.2) et (2.2.4)). On a par
suite dans $\mathcal D$ un diagramme :

\VerdierSemanticAnchor{ch2-2-3-4-f-diagram}
\begin{center}
  \begin{tikzcd}[column sep=4.1em,row sep=2.5em]
    Z \arrow[r,"t"] & X & Y \arrow[l,"s"']
  \end{tikzcd}
  \quad .
\end{center}

\VerdierSemanticAnchor{ch2-2-3-4-f-retraction}
D'après a), iv), le morphisme $s$ admet une rétraction
$s'\colon X\longrightarrow Y$ qui est un élément de
$S_{\mathcal D}(\mathcal B)$ (1.2.6). Par suite, il existe un morphisme
$t'\colon Z\longrightarrow Y$,
$t'\in\widehat{S_{\mathcal D}(\mathcal B)}$.

\VerdierSemanticAnchor{ch2-2-3-4-f-converse}
Réciproquement, s'il existe un tel morphisme, le morphisme
$Q(t')\colon Q(Z)\longrightarrow Q(Y)$ est un isomorphisme.

\VerdierSemanticAnchor{ch2-2-3-4-g}
Il reste à démontrer g). On a, d'après a), v), un isomorphisme
bifonctoriel :

\VerdierSemanticAnchor{ch2-2-3-4-g-q1-adjunction}
\[
  \operatorname{Hom}_{Q^{-1}(\operatorname{Lib}_R(Q))}(Z,\Psi(W))
  \xrightarrow{\sim}
  \operatorname{Hom}_{\operatorname{Lib}_R(Q)}(Q_1(Z),W),
\]
ce qui montre que $\Psi$ est adjoint à droite à $Q_1$.

\VerdierSemanticAnchor{ch2-2-3-4-g-inclusion-adjunction}
On en déduit immédiatement que si on note
$i\colon L_R(Q)\lhook\joinrel\longrightarrow
Q^{-1}(\operatorname{Lib}_R(Q))$ le foncteur d'inclusion, on a un
isomorphisme bifonctoriel :
\[
  \operatorname{Hom}_{Q^{-1}(\operatorname{Lib}_R(Q))}(Z,i(Y))
  \xrightarrow{\sim}
  \operatorname{Hom}_{L_R(Q)}(\Psi Q_1(Z),Y)\, .
\]

\VerdierSemanticAnchor{ch2-2-3-4-g-proof-complete}
Ceci montre que le foncteur $\mathcal R=\Psi Q_1$ est adjoint à gauche au
foncteur d'inclusion et achève la démonstration de la proposition (2.3.3).

\begingroup\itshape
\VerdierSemanticAnchor{ch2-2-3-5}
\medskip
\noindent\textbf{Proposition 2.3.5.}\quad
\VerdierSemanticAnchor{ch2-2-3-5-setup}
Soient $\mathcal D$ une catégorie triangulée, $\mathcal A$ et $\mathcal B$
deux sous-catégories triangulées pleines ($\mathcal B$ strictement pleine),
$S_{\mathcal D}(\mathcal B)$ et
$S_{\mathcal A}(\mathcal A\cap\mathcal B)$ les systèmes multiplicatifs de
$\mathcal D$ et $\mathcal A$ associés respectivement aux catégories
$\mathcal B$ et $\mathcal A\cap\mathcal B$.

\VerdierSemanticAnchor{ch2-2-3-5-a}
\par\noindent a) Les conditions suivantes sont équivalentes :

\VerdierSemanticAnchor{ch2-2-3-5-a-i}
\par\noindent\hspace{1.2em}i) Tout morphisme
$X\xrightarrow{s}X'$, avec $X\in\operatorname{Ob}(\mathcal A)$,
$s\in S_{\mathcal D}(\mathcal B)$, s'insère dans un diagramme commutatif :

\VerdierSemanticAnchor{ch2-2-3-5-a-i-diagram}
\begin{center}
  \begin{tikzcd}[column sep=4.1em,row sep=2.6em]
    X \arrow[r,"s"] \arrow[dr,"{s'}"'] & X' \arrow[d] &
    {s'\in S_{\mathcal A}(\mathcal A\cap\mathcal B)} \\
    & X'' &
  \end{tikzcd}
  \quad .
\end{center}
\endgroup
\endgroup
